\documentclass[11pt]{article}

\usepackage[T1]{fontenc}
\usepackage{lmodern}
\usepackage{amsmath,amssymb,amsthm,mathtools}
\usepackage{microtype}
\usepackage[margin=1.12in]{geometry}
\usepackage{xcolor}
\usepackage[colorlinks=true,linkcolor=blue!55!black,citecolor=blue!55!black,
  urlcolor=blue!55!black]{hyperref}
\usepackage[nameinlink,noabbrev]{cleveref}

\newtheorem{theorem}{Theorem}[section]
\newtheorem{proposition}[theorem]{Proposition}
\newtheorem{corollary}[theorem]{Corollary}
\newtheorem{lemma}[theorem]{Lemma}
\theoremstyle{definition}
\newtheorem{definition}[theorem]{Definition}
\theoremstyle{remark}
\newtheorem{remark}[theorem]{Remark}

\newcommand{\C}{\mathbb C}
\newcommand{\N}{\mathbb N}
\newcommand{\Pp}{\mathcal P}
\newcommand{\A}{\mathcal A}
\newcommand{\K}{\mathcal K}
\newcommand{\Kfin}{\mathcal K_{\mathrm{fin}}}
\newcommand{\Q}{\mathcal Q}
\newcommand{\Rloc}{\mathcal R_{\mathrm{ray}}}
\newcommand{\Span}{\operatorname{span}}
\newcommand{\DA}{\operatorname{DA}}
\newcommand{\HT}{\operatorname{HT}}
\newcommand{\Ann}{\operatorname{Ann}}
\newcommand{\ord}{\operatorname{ord}}

\title{The Uniform-Annihilator Law for Prime-Profile Hypertranscendence:\\
Exact Inverse Limits and Unavoidable Clock-Order Escape}

\author{Michael Haufschild \and GPT 5.6 Sol \and Fable 5 \and Grok 4.5}
\date{}

\begin{document}
\maketitle

\begin{abstract}
Differential algebraicity is witnessed by a finite-order equation, but an infinite intersection
of differential fields can destroy every equation visible at a finite stage.  We give an exact
classification of this phenomenon for prime-power Dirichlet profiles.  Let $\Q$ be the vector
space of complete prime-power profiles modulo profiles carried by finitely many prime rays, let
$T$ be the descended logarithmic clock, and let $\Phi(W)$ be the embedded arithmetic-function
field represented by a source subspace $W\leq\Q$.  We prove that $\Phi$ preserves every nonempty
intersection, not merely finite ones.  If the $W_i$ are $T$-invariant and $f$ is a prime-power
source, then the carrier of $f$ is differentially algebraic over
$\bigcap_i\Phi(W_i)$ if and only if one nonzero finite polynomial $P$ satisfies
$P(T)[f]\in W_i$ for every $i$.  Stagewise differential algebraicity has the weaker quantifier
order $\forall i\,\exists P_i$.  Thus the entire inverse-limit defect is the failure of a uniform
finite clock annihilator.

For descending chains we prove a bounded-order compactness law: if the annihilators at all stages
have a common order bound, then one annihilator survives the limit.  Consequently every genuine
stagewise-algebraic/limit-hypertranscendental transition forces the annihilator order to escape
every finite bound.  The law is sharp.  An explicit source, equal to one on every positive prime
power, has a free clock orbit; its successive orbit-tail fields form a descending chain over which
the source is differentially algebraic at every stage but hypertranscendental over the exact
intersection.  The result isolates a new finite-law obstruction for hypertranscendence arguments
based on limiting prime data.
\end{abstract}

\medskip
\noindent\textbf{Keywords.}
hypertranscendence; differential algebra; Dirichlet series; arithmetic functions; inverse limits;
annihilators; prime-power profiles; Grassmann lattices.

\medskip
\noindent\textbf{2020 Mathematics Subject Classification.}
Primary 12H05; Secondary 11J85, 11M41, 15A03.

\section{Introduction}

Let $L$ be a differential field, $F\subseteq L$ a differential subfield, and $x\in L$.
Differential algebraicity of $x$ over $F$ is a finite statement: some nonzero polynomial in
$x,x',\ldots,x^{(r)}$, with coefficients in $F$, vanishes.  This finiteness makes infinite field
intersections delicate.  From
\[
  x\in\DA\!\left(\bigcap_i F_i\right)
  \quad\Longrightarrow\quad
  x\in\bigcap_i\DA(F_i)
\tag{1.1}
\]
the converse does not follow, because the equation may depend on $i$.  The elementary diagnosis
``the equations are not uniform'' is correct but incomplete: it does not identify a natural class
of fields for which uniformity is the \emph{only} obstruction, nor does it measure the price of
failure.

This paper solves that problem for a source-faithful family of differential fields built from
prime-power profiles.  Three structures interact:
\[
  \boxed{
  \begin{gathered}
  \text{source subspaces }W\leq\Q,
  \qquad \text{embedded fields }\Phi(W),
  \qquad \text{finite clock laws }P(T)[f]\in W.\tag{1.2}
  \end{gathered}}
\]
The quotient $\Q$ removes changes on finitely many prime rays.  The operator $T$ is forced by
logarithmic differentiation.  The field $\Phi(W)$ remembers exactly the source subspace $W$ and,
when $W$ is $T$-invariant, is already a differential field.

Our main result is the following uniform-annihilator law.  The notation is made precise in
\cref{sec:arithmetic,sec:grassmann}.

\begin{theorem}[Uniform-annihilator inverse-limit law]\label{thm:master-intro}
Let $(W_i)_{i\in I}$ be a nonempty family of $T$-invariant subspaces of $\Q$, and let $f$ be
supported on positive prime powers.  Put $q=[f]\in\Q$ and $x_f\in\K$ for its carrier.  Then
\begin{align}
  \bigcap_{i\in I}\Phi(W_i)
    &=\Phi\!\left(\bigcap_{i\in I}W_i\right), \tag{1.3}\\
  x_f\in\DA\!\left(\bigcap_{i\in I}\Phi(W_i)\right)
    &\Longleftrightarrow
      \exists P\in\C[X]\setminus\{0\}\quad
        \forall i\in I,\ P(T)q\in W_i, \tag{1.4}\\
  \forall i\in I,\ x_f\in\DA(\Phi(W_i))
    &\Longleftrightarrow
      \forall i\in I\quad\exists P_i\in\C[X]\setminus\{0\},\quad
        P_i(T)q\in W_i. \tag{1.5}
\end{align}
\end{theorem}

The theorem turns a potentially nonlinear differential-algebra problem into a precise quantifier
switch:
\[
  \exists P\ \forall i
  \qquad\text{versus}\qquad
  \forall i\ \exists P_i.                                      \tag{1.6}
\]
No additional common field element is hidden in the intersection: equation~(1.3) removes that
possibility.  No nonlinear differential relation escapes the clock test: equation~(1.4) removes
that possibility.  What remains is exactly uniform finite annihilation.

For a descending chain the classification has a stronger consequence.

\begin{theorem}[Bounded-order compactness and escape]\label{thm:escape-intro}
Suppose $W_{n+1}\subseteq W_n$ and every $W_n$ is $T$-invariant.  If there is an $r$ such that,
for every $n$, one can choose a nonzero $P_n$ with
\[
  \deg P_n\leq r,
  \qquad P_n(T)q\in W_n,                                      \tag{1.7}
\]
then $x_f$ is differentially algebraic over $\bigcap_n\Phi(W_n)$.  Equivalently, if $x_f$ is
differentially algebraic over every $\Phi(W_n)$ but hypertranscendental over their intersection,
then
\[
  \forall r\in\N\quad\exists n\in\N\quad
  \text{no nonzero annihilator of order at most $r$ works at stage $n$}. \tag{1.8}
\]
\end{theorem}

Thus a differential equation cannot simply evaporate under an inverse limit while its complexity
stays bounded.  Its clock order must escape to infinity.  This is the positive law exposed by the
failure of the converse to~(1.1).

The law is not vacuous.  We construct a source $u$ whose coefficient is one at every positive
prime power.  Its clock orbit is free in $\Q$.  If
\[
  W_N=\Span_{\C}\{T^{N+k}[u]:k\geq0\},                         \tag{1.9}
\]
then $\bigcap_NW_N=0$, the fields $\Phi(W_N)$ have exact intersection $\Phi(0)$, and
\[
  x_u\in\DA(\Phi(W_N))\quad\text{for every $N$},
  \qquad x_u\in\HT(\Phi(0)).                                  \tag{1.10}
\]
This example lies in the actual prime-ray quotient and in its embedded carrier fields.  It is not
the formal example of adjoining a free differential indeterminate and then taking derivative-tail
fields.

Classical differential algebra supplies the language of differential dependence
\cite{Kolchin1973}.  Arithmetic-function rings and formal Dirichlet series supply the ambient
algebra \cite{CashwellEverett1959}; algebraic independence of arithmetic functions and Dirichlet
series has a substantial literature, including \cite{ShapiroSparer1986,Pong2016}.  Recent work on
relative differential closure studies intersections in the very different setting of Hardy fields
\cite{AschenbrennerVanDenDriesVanDerHoeven2025}.  The contribution here is not any of those parent
ingredients separately.  It is the exact conjunction of complete source-field intersection,
nonlinear-to-linear clock reduction, bounded-order compactness, and a strict prime-supported
fracture.

\section{The arithmetic differential field}\label{sec:arithmetic}

Let $\A$ be the ring of complex arithmetic functions $a:\N_{>0}\to\C$, with pointwise addition
and Dirichlet convolution
\[
  (a*b)(n)=\sum_{d\mid n}a(d)b(n/d).
\tag{2.1}
\]
The ring $\A$ is an integral domain; write $\K=\operatorname{Frac}(\A)$ for its fraction field.
Define
\[
  (\delta a)(n)=\mathrm i\log(n)\,a(n).                         \tag{2.2}
\]
Since $\log(dm)=\log d+\log m$, the map $\delta$ is a derivation of $(\A,*)$ and extends uniquely
to $\K$.  Under the formal Dirichlet transform $a\mapsto\sum_{n\geq1}a(n)n^{-s}$, this is the
coefficient action of differentiation in the vertical variable.  No convergence assumption is
needed in what follows; for the analytic theory of general Dirichlet series, see
\cite{DefantSchoolmann2019}.

For each prime $p$, there is also a prime-coordinate derivation detecting the $p$-direction of a
fraction.  Let $\Kfin\subseteq\K$ be the subfield of fractions detected by only finitely many of
these prime derivations.  It contains every alteration carried by finitely many prime rays and is
stable under $\delta$.  It is the coefficient field throughout the paper.

Let $\Pp$ be the set of primes.  A \emph{positive prime-power source} is an arithmetic function
$f$ satisfying
\[
  f(n)\ne0\quad\Longrightarrow\quad n=p^{e+1}
  \quad\text{for some }(p,e)\in\Pp\times\N.                    \tag{2.3}
\]
Its complete profile is
\[
  c_f(p,e)=f(p^{e+1}).                                         \tag{2.4}
\]
Let $\Rloc\leq\C^{\Pp\times\N}$ consist of profiles supported on only finitely many prime rays,
with no restriction along any one ray, and put
\[
  \Q:=\C^{\Pp\times\N}/\Rloc.                                 \tag{2.5}
\]
We write $[f]\in\Q$ for the class of $c_f$.

The logarithmic clock on profiles is
\[
  (Tc)(p,e)=\mathrm i(e+1)\log p\,c(p,e).                      \tag{2.6}
\]
It preserves $\Rloc$ and therefore descends to a linear endomorphism $T:\Q\to\Q$.  If
\[
  f^{\langle k\rangle}(n)=(\mathrm i\log n)^k f(n),            \tag{2.7}
\]
then
\[
  \delta^k x_f=x_{f^{\langle k\rangle}},
  \qquad [f^{\langle k\rangle}]=T^k[f],                        \tag{2.8}
\]
where $x_f$ denotes the image of $f$ in $\K$.

\section{Source subspaces and represented fields}\label{sec:grassmann}

For a subspace $W\leq\Q$, define $\Phi(W)$ to be the intermediate field of $\K/\Kfin$ generated
by the carriers $x_f$ of positive prime-power sources whose classes lie in $W$.  Equivalently, one
may choose a Hamel basis of $W$, choose one positive prime-power representative of each basis
vector, and adjoin only those representatives.  The resulting field is independent of these
choices.

The following finite-level correspondence is the foundation of the inverse-limit theorem.

\begin{proposition}[Prime-profile field representation]\label{prop:representation}
For subspaces $U,W\leq\Q$,
\begin{align}
  \Phi(U)\subseteq\Phi(W)&\Longleftrightarrow U\subseteq W,      \tag{3.1}\\
  \Phi(U+W)&=\Phi(U)\Phi(W),                                    \tag{3.2}\\
  \Phi(U\cap W)&=\Phi(U)\cap\Phi(W).                            \tag{3.3}
\end{align}
Moreover, for every positive prime-power source $f$,
\[
  x_f\in\Phi(W)\Longleftrightarrow [f]\in W.                   \tag{3.4}
\]
If $T(W)\subseteq W$, then $\Phi(W)$ is stable under $\delta$.
\end{proposition}

The crucial strengthening needed here is preservation of arbitrary nonempty meets.

\begin{lemma}[Finite source support]\label{lem:finite-source}
If $z\in\Phi(W)$, then there is a finite-dimensional subspace $U\leq W$ such that
$z\in\Phi(U)$.
\end{lemma}

\begin{proof}
Choose a basis presentation of $\Phi(W)$.  By definition, $\Phi(W)$ is the field generated over
$\Kfin$ by the corresponding carriers.  A rational expression for $z$ uses only finitely many of
those generators.  Let $U$ be the span of their basis classes.  Then $U$ is finite-dimensional,
$U\leq W$, and the same rational expression places $z$ in $\Phi(U)$.
\end{proof}

\begin{lemma}[Finite determination inside a finite source]\label{lem:finite-determination}
Let $U$ be finite-dimensional and let $(W_i)_{i\in I}$ be any family of subspaces of $\Q$.  There
is a finite set $S\subseteq I$ such that
\[
  U\cap\bigcap_{i\in I}W_i=U\cap\bigcap_{i\in S}W_i.           \tag{3.5}
\]
\end{lemma}

\begin{proof}
The subspaces $U\cap\bigcap_{i\in S}W_i$, indexed by finite $S\subseteq I$, form a downward
directed family inside the finite-dimensional space $U$.  Choose one with minimal dimension.
Intersecting it with any further $W_i$ cannot lower the dimension and therefore leaves it
unchanged.  It is consequently contained in every $W_i$, which proves~(3.5).
\end{proof}

\begin{theorem}[Complete nonempty meet law]\label{thm:complete-meet}
For every nonempty family $(W_i)_{i\in I}$ of subspaces of $\Q$,
\[
  \bigcap_{i\in I}\Phi(W_i)=\Phi\!\left(\bigcap_{i\in I}W_i\right). \tag{3.6}
\]
\end{theorem}

\begin{proof}
Monotonicity of $\Phi$ gives the inclusion from right to left.  Conversely, take
$z\in\bigcap_i\Phi(W_i)$ and fix $i_0\in I$.  By \cref{lem:finite-source}, there is a
finite-dimensional $U\leq W_{i_0}$ with $z\in\Phi(U)$.  Apply
\cref{lem:finite-determination} to obtain finite $S\subseteq I$ satisfying~(3.5).  Since
$z\in\Phi(W_i)$ for every $i$, the binary meet law~(3.3), iterated over $S$, yields
\[
  z\in\Phi(U)\cap\bigcap_{i\in S}\Phi(W_i)
   =\Phi\!\left(U\cap\bigcap_{i\in S}W_i\right)
   =\Phi\!\left(U\cap\bigcap_{i\in I}W_i\right).
\tag{3.7}
\]
The last source subspace is contained in $\bigcap_iW_i$, so~(3.1) places $z$ in the right-hand
side of~(3.6).
\end{proof}

The nonempty condition is structural.  The empty meet in the full lattice of intermediate fields
is the ambient field $\K$, whereas $\Phi(\Q)$ need not be all of $\K$.  Thus $\Phi$ is a complete
meet embedding into its represented interval, not an assertion about the empty meet of the
ambient lattice.

\section{The prime-profile Kolchin correspondence}\label{sec:kolchin}

For a $T$-invariant subspace $W\leq\Q$ and a positive prime-power source $f$, define its finite
clock-certificate space by
\[
  \Ann_W(f):=
  \left\{P\in\C[X]:P(T)[f]\in W\right\}.                       \tag{4.1}
\]
This is a complex vector space.  Its nonzero elements are precisely the finite linear clock laws
for $f$ modulo $W$.

We use the usual terminology
\[
  x\in\DA(F)
  \Longleftrightarrow
  (x,\delta x,\delta^2x,\ldots)
  \text{ is algebraically dependent over }F,                   \tag{4.2}
\]
and $x\in\HT(F)$ for the negation.

The special rigidity of the prime-profile sector is that an arbitrary nonlinear differential
dependence is equivalent to a linear clock law.

\begin{lemma}[Exact static prime-profile matroid]\label{lem:static-matroid}
For any family $(g_j)_{j\in J}$ of positive prime-power sources,
\[
  (x_{g_j})_{j\in J}\text{ is algebraically independent over }\Kfin
  \quad\Longleftrightarrow\quad
  ([g_j])_{j\in J}\text{ is linearly independent in }\Q.       \tag{4.3}
\]
\end{lemma}

\begin{proof}
Both sides have finite character, so assume $J$ finite.  If a nontrivial linear combination of the
$[g_j]$ vanishes, the same combination of the profiles is supported on finitely many prime rays.
Its carrier belongs to $\Kfin$, producing a nontrivial linear, hence algebraic, relation among the
$x_{g_j}$ over $\Kfin$.

Conversely, quotient independence says that after deleting any finite set of prime rays, the
remaining ray restrictions of the $g_j$ still have full column rank.  Finite-dimensional row
selection therefore gives distinct fresh prime rays and valuation coordinates on which the
evaluation matrix is nonsingular.  Apply the corresponding prime-coordinate derivations to a
hypothetical polynomial relation of least total degree.  Its Jacobian with respect to those fresh
directions is the product of the nonsingular evaluation matrix with nonzero diagonal prime
factors.  The usual differential Jacobian argument lowers the degree unless the original
polynomial is zero, a contradiction.  Hence the carriers are algebraically independent.  The
arbitrary-family statement follows by applying the finite result to every finite subfamily.
\end{proof}

\begin{lemma}[Relative prime-profile matroid]\label{lem:relative-matroid}
Let $W\leq\Q$ be $T$-invariant, and let $(g_h)_{h\in H}$ be any family of positive prime-power
sources.  The carriers $(x_{g_h})_{h\in H}$ are differentially algebraically independent over
$\Phi(W)$ if and only if the quotient clock family
\[
  \bigl\{T^n[g_h]+W:(h,n)\in H\times\N\bigr\}                  \tag{4.4}
\]
is linearly independent in $\Q/W$.
\end{lemma}

\begin{proof}
Both assertions have finite character, so it suffices to treat a finite packet of sources and
clock orders.  Choose positive prime-power representatives of a basis of $W$.  Since $W$ is
$T$-invariant, their differential field is $\Phi(W)$.

Suppose first that the quotient clock packet is linearly dependent.  Lifting the dependence to
$\Q$ expresses a nontrivial complex linear combination of the new clock jets as an element of
$W$.  Source-field membership~(3.4) turns that element into a coefficient in $\Phi(W)$, and~(2.8)
gives a nontrivial algebraic relation among the differential jets of the $x_{g_h}$.

Conversely, suppose that a polynomial relation among a finite new jet packet holds over
$\Phi(W)$.  Clear its finitely many field denominators.  Every remaining coefficient belongs to a
field generated by finitely many basis representatives of $W$, so the relation is contained in a
finite source presentation.  Select a nonzero coefficient monomial.  By
\cref{lem:static-matroid}, applying the exact static equivalence after adjoining the finite old
source presentation forces a nontrivial
linear relation among the new clock classes modulo its span, hence modulo $W$.  This is the desired
quotient dependence.
\end{proof}

\begin{proposition}[Invariant-Grassmann Kolchin law]\label{prop:kolchin}
Let $W\leq\Q$ be $T$-invariant and let $f$ be supported on positive prime powers.  Then
\[
  x_f\in\DA(\Phi(W))
  \Longleftrightarrow
  \Ann_W(f)\ne\{0\}.                                           \tag{4.5}
\]
Equivalently,
\[
  x_f\in\HT(\Phi(W))
  \Longleftrightarrow
  \{T^n[f]:n\geq0\}\text{ is linearly independent in }\Q/W.   \tag{4.6}
\]
\end{proposition}

\begin{proof}
Apply \cref{lem:relative-matroid} to the singleton family $g_*=f$.  Differential algebraicity is
the negation of differential independence, and linear dependence of the clock orbit modulo $W$
is exactly the existence of a nonzero polynomial $P$ with $P(T)[f]\in W$.
\end{proof}

The proposition is stronger than detecting a special class of linear differential equations.  It
says that within this carrier sector there are no further nonlinear differential-algebraic
relations after the clock quotient has been tested.

\section{The uniform-annihilator law}\label{sec:uniform}

We now prove the main theorem in its exact form.

\begin{theorem}[Uniform-annihilator law]\label{thm:uniform}
Let $(W_i)_{i\in I}$ be a nonempty family of $T$-invariant subspaces of $\Q$, and let $f$ be a
positive prime-power source.  Then
\begin{align}
 x_f\in\DA\!\left(\bigcap_i\Phi(W_i)\right)
 &\Longleftrightarrow
 \bigcap_i\Ann_{W_i}(f)\ne\{0\},                               \tag{5.1}\\
 \forall i,\ x_f\in\DA(\Phi(W_i))
 &\Longleftrightarrow
 \forall i,\ \Ann_{W_i}(f)\ne\{0\}.                            \tag{5.2}
\end{align}
In particular, the failure of differential algebraicity to commute with the intersection is
exactly
\[
  \left(\forall i,\ \Ann_{W_i}(f)\ne\{0\}\right)
  \quad\text{and}\quad
  \bigcap_i\Ann_{W_i}(f)=\{0\}.                                \tag{5.3}
\]
\end{theorem}

\begin{proof}
Arbitrary intersections of $T$-invariant subspaces are $T$-invariant.  By
\cref{thm:complete-meet},
\[
  \bigcap_i\Phi(W_i)=\Phi\!\left(\bigcap_iW_i\right).           \tag{5.4}
\]
The certificate construction commutes with intersections:
\[
  \Ann_{\cap_iW_i}(f)=\bigcap_i\Ann_{W_i}(f),                   \tag{5.5}
\]
because $P(T)[f]$ belongs to $\bigcap_iW_i$ exactly when it belongs to every $W_i$.  Applying
\cref{prop:kolchin} to~(5.4) proves~(5.1).  Applying the same proposition separately at every
stage proves~(5.2), and~(5.3) is their conjunction.
\end{proof}

The theorem rules out two possible sources of inverse-limit pathology.  First, there is no
unrepresented common field sector, by~(5.4).  Second, there is no nonlinear relation invisible to
the linear clock module, by \cref{prop:kolchin}.  The quantifier switch in~(1.6) is therefore a
classification, not merely a sufficient obstruction.

\section{Bounded-order compactness}\label{sec:compactness}

Let $\C[X]_{\leq r}$ be the $(r+1)$-dimensional space of polynomials of degree at most $r$.  Define
\[
  A_n^{(r)}:=\Ann_{W_n}(f)\cap\C[X]_{\leq r}.                   \tag{6.1}
\]
If $W_{n+1}\subseteq W_n$, then $A_{n+1}^{(r)}\subseteq A_n^{(r)}$.

\begin{theorem}[Bounded-order compactness]\label{thm:bounded}
Let $(W_n)_{n\geq0}$ be a descending chain of $T$-invariant subspaces.  If for some $r$
\[
  A_n^{(r)}\ne\{0\}\qquad\text{for every }n,                   \tag{6.2}
\]
then
\[
  \bigcap_n A_n^{(r)}\ne\{0\}                                 \tag{6.3}
\]
and $x_f\in\DA(\bigcap_n\Phi(W_n))$.
\end{theorem}

\begin{proof}
The $A_n^{(r)}$ form a descending chain of subspaces of the finite-dimensional vector space
$\C[X]_{\leq r}$.  Their dimensions stabilize.  Once the dimension stabilizes, every subsequent
inclusion is an equality; hence their intersection is one of the nonzero stabilized spaces.  A
nonzero $P$ in the intersection satisfies $P(T)[f]\in W_n$ for all $n$.  The conclusion follows
from \cref{thm:uniform}.
\end{proof}

\begin{corollary}[Clock-order escape]\label{cor:escape}
Assume $x_f\in\DA(\Phi(W_n))$ for every $n$, but
\[
  x_f\in\HT\!\left(\bigcap_n\Phi(W_n)\right).                  \tag{6.4}
\]
Then for every $r$ there is an $n$ such that $A_n^{(r)}=\{0\}$.
\end{corollary}

\begin{proof}
Otherwise~(6.2) would hold for some $r$, contradicting \cref{thm:bounded}.
\end{proof}

This is a finite-law compactness principle, not a topological compactness argument.  Its source is
the finite-dimensionality of every bounded clock-order sector.  Coefficients may vary without
control and the ambient source spaces may be infinite-dimensional; only the derivative order is
bounded.

\section{A strict prime-supported fracture}\label{sec:fracture}

Define the unit valuation source $u\in\A$ by
\[
  u(p^{e+1})=1\quad(p\in\Pp,\ e\geq0),
  \qquad u(n)=0\quad\text{otherwise}.                           \tag{7.1}
\]
Let $q=[u]\in\Q$.

\begin{lemma}[Free unit-source clock orbit]\label{lem:free-unit}
The sequence $(T^nq)_{n\geq0}$ is linearly independent over $\C$.
\end{lemma}

\begin{proof}
Suppose $P(T)q=0$ for a polynomial $P\in\C[X]$.  Equality to zero in $\Q$ means that, outside a
finite set of prime rays, the corresponding profile vanishes identically.  Choose a prime $p$
outside that set.  On its $e$th valuation coordinate,
\[
  (P(T)q)(p,e)=P\!\left(\mathrm i(e+1)\log p\right).            \tag{7.2}
\]
This vanishes for every $e\geq0$.  The arguments in~(7.2) are pairwise distinct, so $P$ has
infinitely many roots and must be zero.  Hence no nontrivial finite linear combination of the
clock orbit vanishes.
\end{proof}

For $N\geq0$, put
\[
  W_N:=\Span_{\C}\{T^{N+k}q:k\geq0\}.                          \tag{7.3}
\]
Then $T(W_N)\subseteq W_N$ and $W_{N+1}\subseteq W_N$.

\begin{lemma}[Vanishing tail intersection]\label{lem:tail-zero}
\[
  \bigcap_{N\geq0}W_N=0.                                       \tag{7.4}
\]
\end{lemma}

\begin{proof}
An element of $W_0$ is a finite linear combination of the free orbit in
\cref{lem:free-unit}.  If its last possibly nonzero coordinate has index $m$, membership in
$W_{m+1}$ gives a second expression using only orbit vectors of index greater than $m$.  Linear
independence forces every coefficient in both expressions to vanish.
\end{proof}

\begin{theorem}[Strict Kolchin inverse-limit fracture]\label{thm:fracture}
The unit source satisfies
\begin{align}
  \bigcap_{N\geq0}\Phi(W_N)&=\Phi(0),                           \tag{7.5}\\
  x_u&\in\DA(\Phi(W_N))\quad\text{for every }N,                 \tag{7.6}\\
  x_u&\in\HT(\Phi(0)).                                         \tag{7.7}
\end{align}
Consequently
\[
  x_u\in\bigcap_N\DA(\Phi(W_N))
  \setminus\DA\!\left(\bigcap_N\Phi(W_N)\right),              \tag{7.8}
\]
and the annihilator order escapes every finite bound along the tail tower.
\end{theorem}

\begin{proof}
Equation~(7.5) follows from \cref{thm:complete-meet,lem:tail-zero}.  At stage $N$, the monomial
$P_N(X)=X^N$ is a nonzero certificate because $T^Nq\in W_N$; hence~(7.6) follows from
\cref{prop:kolchin}.  At the limit, \cref{lem:free-unit} says that no nonzero polynomial satisfies
$P(T)q=0$, so the same proposition gives~(7.7).  Equation~(7.8) is immediate, and clock-order
escape follows from \cref{cor:escape}.
\end{proof}

The example shows that neither complete field-meet preservation nor exact source reconstruction
forces differential-algebraic closure to preserve infinite meets.  The obstruction lies one level
higher: in the noncompact union of finite clock-order sectors.

\section{Consequences for hypertranscendence strategies}\label{sec:consequences}

The main theorem separates three logically different tasks that are often conflated in limiting
arguments.

\begin{enumerate}
\item \emph{Field convergence.}  One must identify the exact limiting coefficient field.  In the
represented prime-profile sector this is solved by \cref{thm:complete-meet}.

\item \emph{Stagewise differential algebraicity.}  At each stage one may find some finite clock
law.  This is the pointwise condition~(5.2).

\item \emph{Survival at the limit.}  One must produce one law valid at every stage, or at least a
uniform order bound from which \cref{thm:bounded} extracts such a law.
\end{enumerate}

This yields a reusable route constraint.

\begin{corollary}[Finite-law obstruction for prime limits]\label{cor:route}
Consider any argument encoded by a descending chain of invariant prime-profile source spaces.  A
stagewise family of differential equations of uniformly bounded order cannot prove that the
limiting carrier is hypertranscendental over the exact limiting field.  Such a transition requires
unbounded order, an operation not represented by a finite clock polynomial, or departure from the
prime-profile carrier sector.
\end{corollary}

\begin{proof}
The first alternative is \cref{cor:escape}; the remaining alternatives are exactly the ways to
leave its hypotheses.
\end{proof}

For the Riemann zeta function, the classical expansion
\[
  \log\zeta(s)=\sum_{p}\sum_{m\geq1}\frac{1}{m p^{ms}}
  \qquad(\Re s>1)                                               \tag{8.1}
\]
exhibits a complete positive-prime-power profile.  Therefore any proposed zeta argument that
organizes its Euler information into a descending tower of the present represented fields is
subject to \cref{cor:route}.  This does not locate a zero of $\zeta$, and it is not an assertion
that every RH strategy admits such an encoding.  It says something sharper and conditional: once
an argument is inside this encoding, bounded-order finite differential laws cannot disappear at
the limit.  A successful route must explicitly account for the required complexity escape or
identify the structure that lies outside the model.

\section{Novelty boundary and scope}\label{sec:scope}

Several ingredients of the proof are classical and are not claimed as contributions:
\begin{itemize}
\item differential algebraicity is witnessed at finite order;
\item a rational field expression uses finitely many generators;
\item descending chains of subspaces stabilize in finite dimension;
\item a free differential indeterminate gives elementary derivative-tail examples;
\item arithmetic functions under Dirichlet convolution and prime-coordinate methods have an
established theory \cite{CashwellEverett1959,ShapiroSparer1986,Pong2016}.
\end{itemize}

The new theorem surface is the exact prime-profile law formed by their interaction:
\begin{enumerate}
\item the represented arithmetic fields preserve every nonempty source intersection;
\item differential algebraicity over an invariant represented field is exactly finite clock
torsion;
\item differential algebraicity over an inverse limit is exactly a \emph{common} annihilator;
\item bounded annihilator order is compact along descending chains;
\item a concrete prime-supported source realizes strict failure and forced order escape.
\end{enumerate}
Nearby work on Hardy-field intersections concerns relative differential closure in ordered fields
of germs \cite{AschenbrennerVanDenDriesVanDerHoeven2025}; work on Dirichlet-series independence
concerns algebraic or differential independence of specified series
\cite{ShapiroSparer1986,Pong2016}.  These are essential comparison theories, but they do not supply
the source-field meet representation and uniform prime-clock certificate law proved here.

The scope is deliberately precise.  The ambient object is the fraction field of arithmetic
functions and its formal Dirichlet differential structure.  Analytic continuation, boundary
values, growth, and zero distribution are not asserted.  The source must be supported on positive
prime powers, and the source subspaces used as differential bases must be clock-invariant.  These
conditions are what make the classification exact rather than heuristic.

\section{Conclusion}

The inverse-limit problem for prime-profile hypertranscendence has a complete answer:
\[
\boxed{
\begin{gathered}
\bigcap_i\Phi(W_i)=\Phi(\bigcap_iW_i),\\[1mm]
x_f\in\DA(\bigcap_i\Phi(W_i))
\Longleftrightarrow
\exists P\ne0\ \forall i,\ P(T)[f]\in W_i,\\[1mm]
\text{stagewise DA but limit HT}
\Longrightarrow
\text{clock order escapes every finite bound}.
\end{gathered}}
\tag{10.1}
\]
The first line says that the field geometry has no hidden inverse-limit defect.  The second says
that nonlinear differential algebraicity reduces exactly to a common finite law.  The third says
what physical or analytic complexity a strict transition must pay: finite-order coherence cannot
be lost without an unbounded-order cascade.  The explicit unit prime source proves that this phase
transition genuinely occurs inside the arithmetic carrier geometry.

\begin{thebibliography}{99}

\bibitem{AschenbrennerVanDenDriesVanDerHoeven2025}
M. Aschenbrenner, L. van den Dries, and J. van der Hoeven,
\emph{Relative differential closure in Hardy fields},
Monatsh. Math. \textbf{208} (2025), 541--568.

\bibitem{CashwellEverett1959}
E. D. Cashwell and C. J. Everett,
\emph{The ring of number-theoretic functions},
Pacific J. Math. \textbf{9} (1959), 975--985.

\bibitem{DefantSchoolmann2019}
A. Defant and I. Schoolmann,
\emph{Hardy spaces of general Dirichlet series---a survey},
preprint, 2019, arXiv:1902.02073.

\bibitem{Kolchin1973}
E. R. Kolchin,
\emph{Differential Algebra and Algebraic Groups},
Academic Press, New York, 1973.

\bibitem{Pong2016}
W. Y. Pong,
\emph{Applications of differential algebra to algebraic independence of arithmetic functions},
Acta Arith. \textbf{172} (2016), 149--173.

\bibitem{ShapiroSparer1986}
H. N. Shapiro and G. H. Sparer,
\emph{On algebraic independence of Dirichlet series},
Comm. Pure Appl. Math. \textbf{39} (1986), 695--745.

\end{thebibliography}

\end{document}
